% Chapter 6 - The Resulting System and Its Evaluation

\section{The Resulting System and Its Evaluation}\label{sec:evaluation}

This chapter draws the work together and puts it to the test. Across the previous
two chapters the system was built in parts; in use, those parts are one running
product, and it is that product, not the parts, which the thesis set out to
deliver. This chapter therefore presents the resulting system as a single whole
before measuring anything of it, and then measures it against the claim on which
everything rests~(\autoref{subsec:sota_interaction}): that webcam-only,
\acrshort{cpu}-only hand tracking can drive interaction which feels immediate.

Two properties must hold for that claim to stand. The recovered pose must reach
the application quickly enough that the resulting motion is not perceived as
lagging, and the cost of producing it must be modest enough that the whole system
runs on commodity hardware without recourse to a \acrshort{gpu}. The evaluation
tests both --- first the latency of the loop as the user experiences it, then
where within the pipeline that latency is spent, and finally the processor and
memory cost of running the streamer and the Unity application together.

Three instruments are used, each described in turn. The receiver's per-frame
timing log~(\autoref{subsec:receiver}) records software latency over a
long session. A separate, more heavily instrumented profiling pass isolates the
three compute stages of the streamer to show how the per-frame budget is divided
between them. A process-level monitor samples the processor and memory footprint
of both processes over a sustained run. Taken together they establish that the
loop is fast, that its cost is concentrated almost entirely in a single stage, and
that it leaves the host with substantial headroom.



\subsection{The Resulting System}\label{subsec:resulting_system}

Before turning to those measurements, it is worth stating plainly what it is
they characterise. The two preceding chapters each described one half of a single
artefact --- the Python pipeline that turns a frame into a streamed
pose~(\autoref{sec:pipeline}), and the Unity application that turns that stream
into interaction~(\autoref{sec:application}). Run together, the two processes form
one system: a real-time, monocular, \acrshort{cpu}-only hand-tracking and
interaction application, in which a person seated at an ordinary laptop reaches
into a three-dimensional scene with their bare hands and is understood.

In operation it is a single closed loop with the user inside it. The webcam sees
the hands, and two \acrshort{mano} hand avatars appear on screen, posed and placed
to mirror the real hands and tracking every movement of the fingers continuously
thereafter~(\autoref{fig:unity-hands}). From that point the user touches no
controller, wears no headset, and pulls on no glove. A fixed vocabulary of five
gestures~(\autoref{subsec:interaction_layer}) carries every command the system
accepts, each reused from scene to scene in a familiar role: the thumb gestures
move between scenes, wave the companion bear's arm, and step through an undo
history; a closed fist grips a grappling rope in one scene and steers a flying
camera in another; a raised index finger paints; and a menu worn on the wrist ---
summoned by a gesture, turned by the wrist itself, and chosen from by a
gesture~(\autoref{fig:bracelet-menu}) --- selects between modes with no pointer or
on-screen control ever appearing. The hands are at once the controller and,
whenever they take command of an avatar, the thing controlled.

Beneath the interaction lies one unbroken data path, realised end to end. A colour
frame is reduced to 21 landmarks by MediaPipe~(\autoref{subsec:hand_tracking}),
smoothed and anchored in space~(\autoref{subsec:landmark_smoothing}), and resolved
in closed form into a canonical \acrshort{mano} pose~(\autoref{subsec:ik_impl});
that pose, a gesture label, and a placement anchor are serialised as a
\acrshort{json} datagram and sent over \acrshort{udp}~(\autoref{subsec:streaming});
and the receiver reverses each of those steps to drive the on-screen skeleton that
every scene reads from~(\autoref{subsec:receiver}). Each link in that chain was set
out and justified on its own in the earlier chapters; the delivered system is the
chain itself, carrying a hand from in front of the camera to inside the scene,
frame after frame, with nothing along the path that a commodity processor cannot
sustain.

What makes this a system rather than a demonstration is that the one loop serves
four genuinely different applications without alteration. The companion-and-grappling
bear, the deformable face, the editable city, and the painting tool are all built
on the one shared interaction layer --- the same posed skeleton, the same gesture
vocabulary, and the same menu and pointing components~(\autoref{subsec:interaction_layer})
--- yet put the recovered pose to incompatible uses: piloting a physics-driven
avatar, dragging individual mesh vertices, repositioning whole buildings, and
laying down strokes that leave the application as point clouds for surface
reconstruction~(\autoref{subsec:nksr}). A drawing made in the painting tool can
even be carried across into the city and set down by hand as raw
points~(\autoref{subsubsec:scene_city_pcd}), so the scenes are not sealed off from
one another but share a common currency of hand-made geometry. The deliverable, in
short, is a self-contained interactive environment, navigated and operated from end
to end by hand, that asks of its user only the camera already built into a laptop.
It is this system, running as a whole, whose responsiveness and cost the remainder
of the chapter sets out to measure.



\subsection{Experimental Setup}\label{subsec:eval_setup}

All measurements were collected on a laptop equipped with a 13th-generation
Intel Core i5-1334U processor (\SI{1.30}{\GHz} base clock), with no
\acrshort{gpu} involvement at any stage of the pipeline.

The latency figures of~\autoref{subsec:latency_eval} were obtained from the
logging mechanism built into the receiver~(\autoref{subsec:receiver}), which
subtracts each datagram's send timestamp from the time of its consumption and
appends the difference to a file. This measures latency as the application sees
it: the full path from capture through detection, filtering, the analytical solve,
serialisation, transport, and the Unity-side scheduling that precedes consumption.
The per-stage decomposition of~\autoref{subsec:latency_breakdown} comes instead
from a dedicated profiling run of the streamer in which high-resolution timers
bracket the detection, filtering, and solve stages individually; it isolates
compute alone and excludes transport and the Unity side. Because these are two
separate sessions capturing different interaction, their absolute totals are not
directly comparable.
The resource figures of~\autoref{subsec:hardware_util} were recorded by a
process-level monitor that sampled the processor share and resident memory of the
streamer and the Unity process once per second over a five-minute session,
reporting each processor share as a fraction of the machine's total capacity.



\subsection{Software Latency}\label{subsec:latency_eval}

The logging mechanism was enabled over a continuous 13-minute session to collect
a representative sample of per-frame software latency. Over \num{20234} frames the
pipeline achieved a mean latency of \SI{30.68}{\ms}, with a median of
\SI{31.00}{\ms} and a standard deviation of \SI{5.03}{\ms}; the 95th and 99th
percentiles were \SI{37.00}{\ms} and \SI{40.00}{\ms} respectively.

All but the rarest frames therefore fall well below the \SI{60}{\ms}
perceptibility threshold identified in~\autoref{subsec:one_euro_filter}, above
which lag becomes noticeable to users. The maximum observed value of
\SI{428}{\ms} represents an isolated outlier consistent with an OS-level
scheduling interruption or a garbage-collection pause on the Unity side; no
sustained degradation was observed across the session.

\begin{figure}[H]
    \centering
    \captionsetup{justification=centering}
    \includegraphics[width=1.0\linewidth]{img/latency_histogram.pdf}
    \caption{Distribution of software latency over \num{20234} frames.
            Vertical lines mark the mean, median, 95th and 99th
            percentiles. The dashed red line is the
            \SI{60}{\ms} perceptibility threshold~(\autoref{subsec:one_euro_filter}).
            Four outlier frames exceeding \SI{70}{\ms} are omitted from the
            display (maximum observed: \SI{428}{\ms}).}
    \label{fig:latency-hist}
\end{figure}



\subsection{Per-Stage Latency Decomposition}\label{subsec:latency_breakdown}

Knowing that the loop is fast leaves open the question of where its time is spent,
which is the more useful thing to know if it is ever to be made faster. To answer
it, a profiling run instrumented the streamer's three compute stages
independently: the MediaPipe detection of the hand
landmarks~(\autoref{subsec:hand_tracking}), the filtering and three-dimensional
anchoring that follow it~(\autoref{subsec:one_euro_filter}), and the analytical
inverse-kinematics solve that recovers the \acrshort{mano}
pose~(\autoref{subsec:ik_impl}). Over \num{5320} frames the division is stark. The
solve dominates at \SI{89.7}{\percent} of per-frame compute;
detection accounts for \SI{8.1}{\percent} and the filtering and anchoring
stage for only \SI{1.9}{\percent},
with the remaining \SI{0.3}{\percent} attributable to the per-hand bookkeeping
around the solve. The breakdown is shown in~\autoref{fig:latency-pie}.

The practical reading is that the analytical solver is the one component whose
cost matters. Detection and filtering are, by comparison, almost free, which
vindicates the decision to keep the MediaPipe front end and the 1€ filter
deliberately light and to concentrate the pipeline's analytical work in the
pose-recovery stage. Any future effort to reduce latency, or to support more
simultaneous hands, would be spent there and almost nowhere else.

\begin{figure}[H]
    \centering
    \captionsetup{justification=centering}
    \includegraphics[width=1.0\linewidth]{img/micro_benchmark_pie.pdf}
    \caption{Share of per-frame pipeline compute by stage, measured over
            \num{5320} frames in a dedicated profiling run. The analytical
            inverse-kinematics solve accounts for \SI{89.7}{\percent} of compute,
            MediaPipe detection for \SI{8.1}{\percent}, filtering and anchoring for
            \SI{1.9}{\percent}, and per-hand overhead for the remaining
            \SI{0.3}{\percent}.}
    \label{fig:latency-pie}
\end{figure}



\subsection{Hardware Resource Utilisation}\label{subsec:hardware_util}

Latency establishes that the loop is fast; the resource cost establishes that it
is cheap. A process-level monitor recorded the processor share and resident memory
of both the streamer and the Unity application once per second over a five-minute
session (\num{300} samples), with each processor share expressed as a fraction of
the machine's total capacity.

The streamer, which carries the entire vision and kinematics pipeline, averaged
\SI{17.5}{\percent} of total processor capacity (95th percentile \SI{19.2}{\percent},
peak \SI{21.1}{\percent}). The Unity application averaged \SI{14.6}{\percent} (95th percentile
\SI{16.9}{\percent}, peak \SI{18.6}{\percent}). Together the two processes account
for roughly a third of the machine's capacity; the system as a whole averaged
\SI{57.7}{\percent} (95th percentile \SI{73.9}{\percent}, peak \SI{85.8}{\percent}),
the balance being operating-system and background activity, including the monitor
and the camera stack themselves. Memory was modest and, more importantly, flat:
the streamer held a mean resident set of \num{495}\,MB and the Unity process
\num{481}\,MB, neither growing materially across the session. The absence of drift
is the expected consequence of the fire-and-forget design, in which only the
newest datagram is ever retained~(\autoref{subsec:receiver},~\autoref{subsec:udp})
and no per-frame state accumulates. \autoref{fig:hardware-monitor} plots both
quantities over the session.

The complete two-process system therefore runs comfortably inside the budget of a
commodity laptop, using no \acrshort{gpu} and leaving, on average, more than
\SI{40}{\percent} of the processor idle.

\begin{figure}[H]
    \centering
    \captionsetup{justification=centering}
    \includegraphics[width=1.0\linewidth]{img/hardware_monitor_graph.pdf}
    \caption{Processor and resident-memory usage of the streamer and the Unity
            application over a five-minute session (\num{300} one-second samples).
            Processor shares are expressed as a fraction of the machine's total
            capacity. Both processes hold a stable memory footprint across the
            session.}
    \label{fig:hardware-monitor}
\end{figure}



\newpage

\subsection{Discussion}\label{subsec:eval_discussion}

The three measurements converge on a single conclusion. The loop's latency sits
below the threshold at which lag becomes perceptible, so the recovered hands move
with the user rather than visibly behind them; that latency is overwhelmingly the
cost of one stage, the analytical solve, which makes the system's behaviour easy
to reason about and gives any future optimisation a single clear target; and the
whole runs within a fraction of a commodity processor's capacity, with stable
memory and no \acrshort{gpu}.

This is the quantitative form of the argument the application chapter made
qualitatively~(\autoref{subsec:application_summary}). The interactions of the four
scenes --- of a kind that comparable systems deliver only through a head-mounted
display, a controller, or a dedicated peripheral~(\autoref{subsec:sota_interaction})
--- are produced here from a single webcam, on a processor that is left more than
\SI{40}{\percent} idle while doing so. The loop is not merely closed but fast and
cheap, and it is that combination, rather than either property alone, that makes
the accessibility claim of this thesis a measured result rather than an aspiration.
